% figures/chart-regional-bar.tex — Regional cement comparison (real CGD data)
% Source: GET /api/regional/CEMENT_001?source=cgd (snapshot: 2026-05-24)
\begin{figure}[H]
  \centering
  \begin{tikzpicture}
    \begin{axis}[
      width=12cm, height=6.5cm,
      ybar,
      bar width=14pt,
      xlabel={ภูมิภาค},
      ylabel={ราคาปูนซีเมนต์ CGD (บาท/ถุง 50 กก.)},
      symbolic x coords={กลาง,ตะวันออก,ตะวันตก,อีสาน,เหนือ,ใต้},
      xtick=data,
      ymin=160, ymax=225,
      nodes near coords,
      nodes near coords align={vertical},
      every node near coord/.append style={font=\scriptsize},
      grid=both, grid style={dashed,gray!25},
      tick label style={font=\footnotesize},
      label style={font=\footnotesize},
    ]
      % Real CGD regional data (1 province per region as deployed)
      \addplot[fill=orange!50,draw=orange!70!black] coordinates {
        (กลาง,175) (ตะวันออก,178) (ตะวันตก,180) (อีสาน,188) (เหนือ,195) (ใต้,215)
      };
    \end{axis}
  \end{tikzpicture}
  \caption{ราคาปูนซีเมนต์ Type I (50 กก.) ตามภูมิภาค จากแหล่งข้อมูล CGD (ดึงจาก \texttt{/api/regional/CEMENT\_001?source=cgd} ของระบบเมื่อ 24 พ.ค.\ 2569). ราคาในภาคใต้สูงกว่าภาคกลาง 22.9\% (215 vs 175 บาท) สอดคล้องกับต้นทุนการขนส่งระยะไกลจากโรงผลิตในภาคกลางและภาคตะวันออก. ข้อมูลในภาพนี้ใช้จังหวัดอ้างอิง 1 จังหวัดต่อภูมิภาค (กรุงเทพฯ, ชลบุรี, ราชบุรี, ขอนแก่น, เชียงใหม่, สงขลา) ตามขอบเขตข้อมูลที่ระบบได้รวบรวมในช่วงทดลอง.}
  \label{fig:chart-regional}
\end{figure}
